% Two authorization mechanisms on one service, and nothing between them.
% Referenced from chapter 15. Bare tikzpicture; the figure environment,
% caption and label live at the call site.
%
% The point of the picture is the vertical divider. It runs the whole height
% of the figure and breaks only once, exactly at the composer, which is the
% only place the two attributes ever appear together. Nothing crosses from
% one side to the other anywhere else: the left lane is checked and consumed
% entirely inside Accounts, the right lane is never read by Accounts at all,
% and the composer is what decides which of the two the router ever hears
% about.
\begin{tikzpicture}[
    font=\small,
    tp/.style={draw, rounded corners=2pt, minimum width=42mm,
               minimum height=15mm, align=center, font=\scriptsize},
    dirbox/.style={draw, rounded corners=2pt, minimum width=40mm,
                   minimum height=13mm, align=center, font=\scriptsize},
    composer/.style={draw, rounded corners=2pt, minimum width=46mm,
                     minimum height=13mm, align=center, font=\scriptsize,
                     fill=black!8},
    dropped/.style={draw, dashed, rounded corners=2pt, minimum width=40mm,
                    minimum height=20mm, align=center, font=\scriptsize,
                    gray},
    kept/.style={draw, rounded corners=2pt, minimum width=46mm,
                minimum height=24mm, align=center, font=\scriptsize,
                fill=black!8},
    flow/.style={-{Stealth[length=2mm]}},
    divider/.style={gray, dashed},
    note/.style={font=\scriptsize, gray, align=center, inner sep=1pt},
    closing/.style={font=\scriptsize\itshape, align=center}
  ]

  % -- one subgraph, two fields -----------------------------------------------
  \node[note, gray, anchor=south, font=\small] at (6.5,1.65) {\code{accounts}};
  \draw[rounded corners=2pt] (0.3,-2.4) rectangle (12.7,1.5);

  \node[tp] (f1) at (3.2,0.5)
    {\code{Query.customerById}\\\code{[Authorize]}\\\code{HotChocolate.Authorization}};
  \node[tp] (f2) at (9.8,0.5)
    {\code{Customer.email}\\\code{[RequiresScopes]}\\\code{HotChocolate.ApolloFederation}};

  \node[note, anchor=north] at (3.2,-0.85)
    {evaluated here,\\against the principal\\built from the token};
  \node[note, anchor=north] at (9.8,-0.85) {Accounts never reads it};

  % -- both compile to SDL, published either way -------------------------------
  \node[dirbox] (d1) at (3.2,-3.5) {\code{@authorize}\\on \code{customerById}};
  \node[dirbox] (d2) at (9.8,-3.5)
    {\code{@requiresScopes}\\\code{(scopes: [["read:pii"]])}};

  \draw[flow] (f1.south) -- (d1.north);
  \draw[flow] (f2.south) -- (d2.north);

  \node[note, anchor=east] at (2.6,-2.6) {compiles to;\\published in\\the SDL};
  \node[note, anchor=west] at (10.4,-2.6) {compiles to;\\published in\\the SDL};

  % -- the only place the two lanes meet ---------------------------------------
  \node[composer] (comp) at (6.5,-5.6)
    {\code{wgc router compose}\\the composer};

  \draw[flow] (d1) -- (comp);
  \draw[flow] (d2) -- (comp);

  \draw[divider] (6.5,1.75) -- (6.5,-5.0);
  \draw[divider] (6.5,-6.2) -- (6.5,-9.0);

  % -- and the two fates ------------------------------------------------------
  \node[dropped] (drop) at (3.2,-8.0)
    {\textbf{silently dropped}\\by the composer\\reaches the router\\as nothing};
  \node[kept] (keep) at (9.8,-8.0)
    {\textbf{survives composition}\\\code{fieldConfigurations}\\entry in the
     router's\\execution config;\\the router refuses};

  \draw[flow] (6.5,-6.2) -- (6.5,-6.6) -- (3.2,-6.6) -- (drop.north);
  \draw[flow] (6.5,-6.2) -- (6.5,-6.6) -- (9.8,-6.6) -- (keep.north);

  \draw[gray, thick] (1.8,-9.1) -- (4.6,-9.1);

  \node[closing, text width=124mm] at (6.5,-10.0)
    {Nothing on the left produces the directive on the right, and nothing on
     the right knows the left exists. The two attributes share one service and
     meet, once, inside the composer, which keeps one and throws the other
     away.};

\end{tikzpicture}
